\subsection{赋值的例子}

§ 1 第 4 节给出的赋值环的例子为我们提供下面例 1 至例 4：

例 (1) 有限域 F 上的每个赋值都是非真的，因为 \(F^{*}\) 的每个元素都是单位根。

例 (2) 若 K 是域 \(K'\) 的子域，则 \(K'\) 上的赋值在 K 上的限制是 K 上的赋值。

例 (3) 设 \(k\) 是域，\(K = k((T))\)。把每个非零形式幂级数映为它的阶（《代数》第四章 § 5 第 7 节）的映射 \(v\) 是 \(K\) 上的赋值，其阶群为 \(Z\)，环为 \(k[[T]]\)。相伴的位是典范同态 \(f: k[[T]] \to k\) 向 \(k((T))\) 的扩张，其中令 \(f(u) = \infty\)（当 \(u \notin k[[T]]\) 时）。

例 (4) 设 A 是主理想整环，K 是它的分式域，p 是 A 的一个极端元。对 \(x \in K^{*}\)，用 \(v_{p}(x)\) 表示 p 在 x 的分解为极端元素的分解式中的指数（《代数》第七章 §1 第 3 节定理 2）；立即可见 \(v_{p}\) 是一个赋值，其阶群为 Z，环为 \(A_{Ap}\)。由 §1 第 4 节命题 3，我们如此得到（在相差一个等价的意义下）K 上的所有非真的且在 A 上为正的赋值。取 A = Z，我们重新得到 Q 上的 p-进赋值（《一般拓扑学》第九章 §3 第 2 节）；这些赋值是（在相差一个等价的意义下）Q 上仅有的非真赋值（§1 第 4 节命题 3 的推论 1）。取 A = k{[}X{]}（其中 k 是域），k(X) 上的限制在 k 上为非真的那些非真赋值是（在相差一个等价的意义下）：一方面是诸赋值 \(v_{P}\)，其中 P 取遍 k{[}X{]} 的首一不可约多项式的集合；另一方面是由

\[
v (\mathrm{P} / \mathrm{Q}) = \deg (\mathrm{Q}) - \deg (\mathrm{P})
\]

（其中 \(P \in k[X]\)、\(Q \in k[X]\)）所定义的赋值 v（§ 1 第 4 节命题 3 的推论 2）；所有这些赋值显然都以 Z 为阶群，其剩余域是 k 的单生成代数扩张（《代数》第五章 § 3 第 1 节）。

例 (5) 映射 \(P(X, Y) \mapsto P(T, e^{\mathrm{T}})\)（从 \(\mathbf{C}[X, Y]\) 到 \(\mathbf{C}((T))\)）是单射（《实变函数》第四章 §2 命题 9），因此可以扩张为从 \(\mathbf{C}(X, Y)\) 到 \(\mathbf{C}((T))\) 的一个子域上的同构。例 3 中定义的 \(\mathbf{C}((T))\) 上的赋值在这个子域上的限制定义了 \(\mathbf{C}(X, Y)\) 上的一个赋值，它在 C 上非真，其阶群为 Z，剩余域为 C。

第 2 节命题 4 使我们能够构造阶群与剩余域都给定的赋值：

例 (6) 设 \(\Gamma\) 是全序群，\(k\) 是域。设 \(\Gamma_{+}\) 是 \(\Gamma\) 的正元素所成的幺半群，C 是 \(\Gamma_{+}\) 在 \(k\) 上的代数。按定义，C 有一组基 \((x_{\alpha})_{\alpha \in \Gamma_{+}}\)（在 \(k\) 上），其乘法表为 \(x_{\alpha}x_{\beta} = x_{\alpha +\beta}\)。若 \(x = \sum_{\alpha}a_{\alpha}x_{\alpha}\) 是 C 的非零元素，记 \(v(x) = \inf_{a_{\alpha}\neq 0}(\alpha)\)，并记 \(v(0) = +\infty\)；立即可验证映射 \(v\)（从 C 到 \(\Gamma_{\infty}\)）满足第 1 节的条件 \((\mathrm{VL}_{\mathrm{I}})\) 与 \((\mathrm{VL}_{\mathrm{II}})\)，且 C 是整环。设 K 是 C 的分式域，\(w\) 是 K 上延拓 \(v\) 的赋值（第 2 节命题 4）。由于 \(\Gamma\) 的每个元素都是两个正元素之差，\(w\) 以 \(\Gamma\) 为阶群。设 A 是 \(w\) 的环，m 是它的极大理想；我们将证明 A 是 m 与 k（等同于 k.1）的直和，这将证明 \(w\) 的剩余域同构于 k。显然 \(m \cap k = (0)\)。另一方面，用 p 表示 C 的由诸 \(x_{\alpha}\)（其中 \(\alpha > 0\)）生成的理想，K 中每个赋值为 0 的元素 x 都可写成 \((a + y)/(b + z)\) 的形式，其中 \(a \in k^{*}\)、\(b \in k^{*}\)、\(y \in p\)、\(z \in p\)；则

\[
x = a b ^ {- 1} + (b y - a z) b ^ {- 1} (b + z) ^ {- 1}
\]

从而 \(w(x - ab^{-1}) > 0\)，且 \(x \equiv ab^{-1} \pmod{m}\)；这就表明了我们的断言。

若 \(\Gamma = \mathbf{Z} \times \mathbf{Z}\)，则 \(\mathbf{K} = k(\mathbf{X}, \mathbf{Y})\)，于是上述构造给出 \(k(\mathbf{X}, \mathbf{Y})\) 上的一批在 \(k\) 上非真的赋值，其阶群为 \(\mathbf{Z} \times \mathbf{Z}\)，剩余域为 \(k\)。这些赋值依赖于 \(\mathbf{Z} \times \mathbf{Z}\) 上所选的序结构。例如，可以给 \(\mathbf{Z} \times \mathbf{Z}\) 赋以字典序。或者，对无理数 \(\alpha\)，\(\mathbf{Z} \times \mathbf{Z}\) 可等同于 \(\mathbf{R}\) 的一个子群（借助同态 \((m, n) \mapsto m + n\alpha\)，由于 \(\alpha\) 是无理数，这个同态是单射），并赋以由 \(\mathbf{R}\) 上的序诱导的序。

利用第 2 节命题 4 的赋值的其他构造将在 § 10 中描述。

